\section{Sub-Gaussian labels with the same learner prior}
\label{sec:subgaussian-labels}

The raw-output oracle inequalities also allow unbounded labels with a
sub-Gaussian tail. The Gaussian parameter prior remains exactly the
same. Only the likelihood calibration acquires logarithmic factors.
The extension here concerns raw outputs and their ordinary posterior
mean, not the earlier convention clipping predictions to $[-1,1]$.

First suppose $|Y|\leq T$ for a known $T\geq1$. In the proof only,
normalize the prediction observable and label by $T$: use $z_w/T$
and $Y/T$, with $\lambda'=\lambda/T^2$. Write
$S/T=((X_i,Y_i/T))_{i=1}^n$; the inputs are not rescaled.
The likelihood identity is
\[
 \exp\left[-\frac{\omega\widehat R_S(z_w)}{2\lambda}\right]
 =\exp\left[-\frac{\omega\widehat R_{S/T}(z_w/T)}{2\lambda'}\right].
\]
This identity does not rescale any Gaussian weight or change the
learner's raw prediction. The prior moments, amplitudes, and Lipschitz
cover bounds for $z_w/T$ are no larger than those for $z_w$.
A teacher observable $g/T$ is represented by the same original teacher
parameters in this normalized observation map, so its energy cost
remains $E$. A raw-output tube of radius
\[
 \rho=\frac\lambda{4(U+T)}
\]
has normalized radius
$\rho/T=\lambda'/[4(U/T+1)]$.
Thus repeating the proof, rather than changing the model's prior,
and then multiplying the risk inequality by $T^2$, gives
\begin{align}
 \int R(z_w)d\Pi
 \leq\;&R(g)+2\lambda\mathcal A_g(\rho)+C_{d,H,A}\lambda\notag\\
 &+(U+T)\sqrt{\frac{2R(g)\log(2/\delta)}n}
       +\frac{2(U+T)^2\log(2/\delta)}{3n},
 \label{eq:bounded-large-label-oracle}
\end{align}
under the sufficient condition
\begin{equation}
 n\geq C_{d,H,A}T^{2p}\lambda^{-p}
           \log^6\!\frac{2nT^2}{\lambda\delta},
 \qquad p=\frac{(H+1)(d+2)}2,
 \quad 0<\lambda\leq1.
 \label{eq:bounded-large-label-budget}
\end{equation}
The stated range of $\lambda$ ensures $\rho\leq1$, as required by
the teacher-cylinder lemma. As before, the sample event is common to
all learner widths above the teacher width and all $0<\tau\leq1$.

\begin{corollary}[A raw oracle with sub-Gaussian labels]
\label{cor:subgaussian-labels}
Replace the bounded-label assumption by the marginal tail condition
\begin{equation}
 \Pr\{|Y|>t\}\leq2e^{-t^2/(2s^2)}\qquad(t\geq0)
 \label{eq:subgaussian-label-tail}
\end{equation}
for a known $s>0$. Let $n\geq2$, $0<\delta<1$, and put
\[
 T=\max\left\{1,s\sqrt{8\log(8n/\delta)}\right\},
 \qquad u=2s\left(\frac\delta{8n}\right)^2.
\]
Fix a bounded-weight teacher $g$, with $|g|\leq U$, and let
$\rho=\lambda/[4(U+T)]$, where $0<\lambda\leq1$.
Suppose \eqref{eq:bounded-large-label-budget} holds with confidence
parameter $\delta/2$. Define
\begin{align*}
 e={}&2\lambda\mathcal A_g(\rho)+C_{d,H,A}\lambda\\
 &+(U+T)(\sqrt{R(g)}+u)\sqrt{\frac{2\log(4/\delta)}n}
          +\frac{2(U+T)^2\log(4/\delta)}{3n}.
\end{align*}
With probability at least $1-\delta$, simultaneously for all $N\geq m$
and all $0<\tau\leq1$, the posterior trained on the original,
unbounded labels satisfies
\begin{equation}
 \int R(z_w)d\Pi
 \leq R(g)+e+4u\sqrt{R(g)+e}+4u^2.
 \label{eq:subgaussian-label-oracle}
\end{equation}
The ordinary raw posterior mean obeys the same bound.
The Gaussian prior and likelihood form are unchanged.
At fixed teacher resources and fixed $s,d,H,A,\delta$, the sufficient
excess rate remains $\widetilde O(n^{-1/p})$.
\end{corollary}
\begin{proof}
Write $Y_T=\operatorname{clip}_{[-T,T]}Y$, and let $R_T$ denote
population square loss for $(X,Y_T)$. Tail integration gives
\begin{align*}
 \|Y-Y_T\|_2^2
 &=\int_0^\infty2r\Pr\{|Y|>T+r\}\,dr\\
 &\leq4e^{-T^2/(2s^2)}
       \int_0^\infty r e^{-r^2/(2s^2)}\,dr
 =4s^2e^{-T^2/(2s^2)}\leq u^2.
\end{align*}
A union bound gives
\[
 \Pr\{\text{some training label differs from its clipping}\}
 \leq2n e^{-T^2/(2s^2)}
 \leq2n(\delta/(8n))^4<\delta/2.
\]
On the complementary event, training on the original and clipped
samples gives exactly the same posterior for every width and scale.
Apply \eqref{eq:bounded-large-label-oracle} to the IID clipped sample
with failure probability $\delta/2$. Minkowski gives
$\sqrt{R_T(g)}\leq\sqrt{R(g)}+u$, so the resulting posterior bound is
$\int R_T(z_w)d\Pi\leq R_T(g)+e$.

For each fixed observed sample, Minkowski on the product of the
posterior law and the fresh-test distribution gives
\[
 \sqrt{\int R(z_w)d\Pi}
 \leq\sqrt{\int R_T(z_w)d\Pi}+u
 \leq\sqrt{R(g)+e}+2u.
\]
The last inequality also uses
$\sqrt{R_T(g)+e}\leq\sqrt{R(g)+e}+u$.
Squaring proves \eqref{eq:subgaussian-label-oracle}, and Jensen gives
the ordinary-mean statement. The two exceptional events have total
probability at most $\delta$ and are common to all asserted learners.
Since $T$ grows only as a square root of a logarithm, its polynomial
factor in the sample condition changes logarithms but not the power
of $n$. The bound $u=O(n^{-2})$ makes the population clipping error
smaller still.
\end{proof}

The same normalization extends the sharper tanh theorem. Its sample
condition is multiplied by $T^{2s_{d,H}}$, with $\lambda$ and
$\epsilon$ inside the logarithm replaced by $\lambda/T^2$ and
$\epsilon/T^2$. Its raw additive term remains $C_{d,H}\epsilon$ after
multiplication by $T^2$. The same clipping and population-risk argument
therefore preserves its fixed-teacher rate and its growing-class
exponents, up to logarithms. The generic raw growing-class bounds
are preserved for the same reason.

Condition~\eqref{eq:subgaussian-label-tail} includes Gaussian noise
around an arbitrary uniformly bounded regression function, with a
sufficiently large $s$. If the noise has conditional mean zero, its
irreducible variance cancels when comparing two population squared
risks; no correct Gaussian likelihood model is needed for the oracle
argument.
